\documentclass[11pt,a4paper]{article}

% --- Paquetes Esenciales ---
\usepackage[utf8]{inputenc}
\usepackage[spanish]{babel}
\usepackage[T1]{fontenc}
\usepackage{geometry}
\geometry{top=2.5cm, bottom=2.5cm, left=2cm, right=2cm}
\usepackage{xcolor}
\usepackage{titlesec}
\usepackage{enumitem}
\usepackage{tcolorbox}
\usepackage{hyperref}
\usepackage{graphicx}
\usepackage{listings}
\usepackage{fancyhdr}
\usepackage{multicol}

% --- Definición de Colores (Onyx & Gold Theme) ---
\definecolor{onyx}{HTML}{1A1A1A}
\definecolor{gold}{HTML}{F59E0B}
\definecolor{pastelyellow}{HTML}{FEF3C7}
\definecolor{softgray}{HTML}{F3F4F6}
\definecolor{darkgray}{HTML}{374151}

% --- Configuración de Estilos ---
\pagestyle{fancy}
\fancyhf{}
\lhead{\textcolor{gold}{\textbf{GANESHA}} | Documentación Técnica}
\rhead{Integrantes del Proyecto}
\cfoot{\thepage}

\titleformat{\section}{\color{onyx}\normalfont\Large\bfseries}{\thesection}{1em}{}[{\titlerule[0.8pt]}]
\titleformat{\subsection}{\color{darkgray}\normalfont\large\bfseries}{\thesubsection}{1em}{}

\hypersetup{
    colorlinks=true,
    linkcolor=gold,
    urlcolor=blue
}

% --- Estilo de Cajas ---
\newtcolorbox{infobox}[1]{
    colback=pastelyellow!30,
    colframe=gold,
    fonttitle=\bfseries,
    title=#1,
    arc=4pt,
    outer arc=4pt
}

\newtcolorbox{codebox}[1]{
    colback=softgray,
    colframe=darkgray,
    fonttitle=\bfseries,
    title=#1,
    arc=2pt
}

% --- Título del Documento ---
\title{
    \vspace{-2cm}
    \begin{center}
        \textcolor{gold}{\Huge \textbf{GANESHA}}\\[0.5em]
        \textcolor{onyx}{\Large \textit{Sistema ERP de Inteligencia Financiera Local}}\\[1cm]
        \textbf{Documento de Arquitectura y Operaciones}\\[0.5em]
        \small{Versión 2.5.0-Enterprise}
    \end{center}
}
\author{Equipo de Desarrollo GANESHA}
\date{\today}

\begin{document}

\maketitle

\begin{abstract}
GANESHA es una plataforma ERP (Enterprise Resource Planning) de próxima generación diseñada específicamente para PYMES en el mercado nicaragüense. Su arquitectura prioriza la soberanía de datos mediante un enfoque \textit{Local-First}, integrando Inteligencia Artificial descentralizada para la toma de decisiones contables sin comprometer la privacidad del usuario.
\end{abstract}

\tableofcontents
\newpage

% --- SECCIÓN 1: MISIÓN Y VISIÓN ---
\section{Visión del Proyecto}
GANESHA no es solo un software contable; es un asesor financiero activo. El objetivo principal es democratizar el acceso a herramientas de alta tecnología (IA y ERP robusto) eliminando la dependencia de servicios en la nube costosos e intrusivos.

\begin{infobox}{Pilar Fundamental: Soberanía de Datos}
Basado en la \textbf{Ley 46 de Nicaragua}, GANESHA asegura que el 100\% de la información contable permanezca en la infraestructura local del cliente, utilizando motores de base de datos embebidos y procesamiento de IA local.
\end{infobox}

% --- SECCIÓN 2: ARQUITECTURA TECNOLÓGICA ---
\section{Arquitectura Tecnológica (Tech Stack)}
El sistema utiliza un stack moderno optimizado para rendimiento y portabilidad:

\begin{itemize}[label=\textcolor{gold}{\textbullet}]
    \item \textbf{Front-end}: Next.js 15 (React 19) con App Router.
    \item \textbf{Diseño}: Tailwind CSS 4 con sistema de temas dinámicos (Onyx, Vintage, Pastel).
    \item \textbf{Back-end}: API Routes de Next.js (Edge-ready) con validación estricta vía Zod.
    \item \textbf{Base de Datos}: SQLite gestionado a través de Prisma ORM para máxima velocidad en entornos locales.
    \item \textbf{IA Local}: Integración con Ollama (Llama 3) para procesamiento de lenguaje natural contable.
\end{itemize}

% --- SECCIÓN 3: MÓDULOS DEL SISTEMA ---
\section{Módulos Funcionales (100\% Cobertura)}

\subsection{Núcleo Contable}
\begin{itemize}
    \item \textbf{Catálogo de Cuentas}: Estructura jerárquica de cuentas (Activos, Pasivos, Capital, Ingresos, Egresos).
    \item \textbf{Libro Diario (Pólizas)}: Registro de transacciones con validación de \textbf{Partida Doble} en tiempo real.
    \item \textbf{Centros de Costo}: Segmentación departamental de gastos e ingresos.
    \item \textbf{Periodos Fiscales}: Gestión de cierres mensuales y anuales con bloqueo de ediciones.
\end{itemize}

\subsection{Tesorería y Operaciones}
\begin{itemize}
    \item \textbf{Facturación}: Gestión de facturas de compra y venta con seguimiento de estados.
    \item \textbf{Bancos y Cuentas}: Conciliación bancaria y flujo de caja.
    \item \textbf{Terceros}: Directorio unificado de Clientes, Proveedores y Empleados.
    \item \textbf{Activos Fijos}: Registro y control de depreciación de bienes.
\end{itemize}

\subsection{Inteligencia y Reportes}
\begin{itemize}
    \item \textbf{Asesor IA}: Chat interactivo que analiza el balance y ofrece consejos financieros.
    \item \textbf{Reportes Financieros}: Balance General, Estado de Resultados y Balanza de Comprobación (Exportables a PDF/Excel).
    \item \textbf{Dashboard de KPIs}: Visualización gráfica de salud financiera y liquidez.
\end{itemize}

% --- SECCIÓN 4: SEGURIDAD Y ROLES ---
\section{Seguridad y Control de Acceso (RBAC)}
El sistema implementa un control de acceso basado en roles (RBAC) para garantizar que cada usuario vea solo lo necesario:

\begin{table}[h]
\centering
\begin{tabular}{|l|p{10cm}|}
\hline
\rowcolor{gold!20} \textbf{Rol} & \textbf{Permisos y Alcance} \\ \hline
\textbf{ADMIN} & Acceso total, gestión de empleados, configuración de empresa y auditoría. \\ \hline
\textbf{ACCOUNTANT} & Gestión completa de pólizas, reportes e impuestos. No crea usuarios. \\ \hline
\textbf{MANAGER} & Visualización de KPIs y reportes para toma de decisiones. Sin permisos de edición. \\ \hline
\textbf{VIEWER} & Acceso de \textbf{Solo Lectura}. No puede crear, editar ni eliminar ningún dato. \\ \hline
\end{tabular}
\end{table}

\begin{infobox}{Seguridad de Contraseñas}
Las credenciales se protegen mediante el algoritmo \textbf{Bcrypt} con un factor de costo de 10, asegurando resistencia contra ataques de fuerza bruta.
\end{infobox}

% --- SECCIÓN 5: AUDITORÍA Y TRAZABILIDAD ---
\section{Bitácora de Auditoría}
GANESHA registra cada acción crítica (Creación de pólizas, edición de usuarios, exportación de datos). Cada registro incluye:
\begin{enumerate}
    \item \textbf{Usuario}: Quién realizó la acción.
    \item \textbf{Acción}: Qué se hizo (CREATE, UPDATE, DELETE).
    \item \textbf{Entidad}: Sobre qué objeto (Póliza \#101, Factura \#FC-001).
    \item \textbf{Timestamp}: Fecha y hora exacta con precisión de milisegundos.
\end{enumerate}

% --- SECCIÓN 6: CONCLUSIÓN ---
\section{Conclusión para Integrantes}
GANESHA representa la vanguardia tecnológica en software empresarial local. Su diseño modular permite una escalabilidad infinita y su enfoque en la privacidad lo posiciona como la opción más segura para empresas que valoran su soberanía financiera.

\begin{center}
\vspace{1cm}
\textcolor{gold}{\rule{5cm}{1pt}}\\
\textbf{GANESHA - Intelligence}\\[0.2em]
"El futuro de la contabilidad es local, inteligente y soberano."
\end{center}

\end{document}
